\chapter{Differential Persuasion and Disclosure}\label{ch:persuasion}

The gravest misuse of selective cinema does not require restricting anyone's glasses. It requires only that different people be shown different things while believing they saw the same film. This chapter examines that possibility and proposes the disclosure that would prevent it.

\section{The hazard}

A multiplexed film could show different products to different viewers, different political messages, different characterizations of the same public figure, or different moral conclusions to the same story. Each viewer would see a coherent film. The viewers might discuss it afterward, each assuming their film was the film. Critics would review one realization and believe they had reviewed the work. Someone who objected to a message would find that others in the same screening had not seen it.

Differential messaging is already common in advertising and political communication, which target different groups with different content. What selective cinema would add is its appearance of commonality. A film screened in a public room has always functioned as evidence that the audience received one text. That is the assumption that makes public criticism possible: there is an object, and anyone can go and see it. A selective film with undisclosed variation would keep the appearance of a common object while removing the object. Criticism would lose its target, because no accessible version contains all the variations and nobody knows which they saw.

The danger is sharpest when variation is combined with the adaptive selection of \Cref{ch:choice}: a system that observes each viewer and then chooses what they see. That combination is targeted persuasion presented as a public screening, and it is the use this book most wants to see excluded.

\section{Disclosure of existence}

The first requirement, already stated in \Cref{ch:after-screening}, is that the existence of variation be disclosed. A viewer entering a selective screening must be told that the audience may be seeing different realizations. This single requirement removes the worst of the hazard, because it removes the false assumption of commonality. A viewer who knows that others may be seeing something else will not take their own film for the film.

\section{A variant registry}

Disclosure of existence is necessary and not sufficient. A viewer who knows variation exists still needs to know what it consists of. The second requirement is a \emph{variant registry}: a public record, for each selective work, of what realizations exist, what dimensions they vary along, what each contains, and who made each. The registry is the public version of the variant matrix of \Cref{ch:matrix}, or of the parts of it that concern what viewers see.

The registry makes criticism possible again. A critic can review the family, not a member of it. A viewer who objects to something in their realization can find out whether others saw it. A researcher can compare what was shown to whom.

\section{Manifests and verification}

A registry says what a work contains. It does not prove that what was screened matches it. A distributor could register one set of realizations and project another.

The third requirement is a \emph{manifest}: a record of what each stream contained as exhibited, with a cryptographic fingerprint of each stream's content, published with the registry. Anyone with access to the exhibited material, a regulator, a critic, an archivist, can check that the fingerprints match. A registered realization that differs from what was shown will not match its fingerprint.

The manifest also answers a question each viewer may reasonably ask: which realization did I see? In the design of \Cref{ch:phase}, the glasses know only their own dial position and the timing signal. A viewer can read their own phase from the glasses and look it up in the public manifest for that screening. Nothing about the viewer is recorded or transmitted, and no one else learns which realization they chose. The viewer learns what they saw without the theater learning who they are.

\section{Reference screenings and audit}

The fourth requirement is that every realization be available for inspection. A selective work should be shown, at least in reference screenings, in a mode in which any realization can be watched in full, including the composite. Archives should hold every realization, as \Cref{ch:which-film} argues. A work whose realizations cannot all be seen by someone accountable is a work whose content is, in part, secret, and a secret shown in public to different people is the hazard itself.

\section{Sponsors}

Commercial pressure deserves a specific word. A sponsor could pay to have its product appear in one realization, or to have a realization made for an audience it wants to reach. There is nothing inherently wrong with product placement, but its legitimacy in ordinary film rests on its visibility: everyone sees the same placements, and anyone can point them out. Placements that vary by realization, invisibly, are a different practice. They should be listed in the registry like any other variation, with the sponsor named, and no realization should exist primarily to carry a message that a sponsor paid for and the audience was not told about.

\section{What disclosure preserves}

The requirements of this chapter are not burdens on the form. They are what allows it to keep the promise of \Cref{ch:synchronized-difference}. Synchronized difference is valuable because it makes divergence visible, co-present, and discussable, the opposite of the hidden divergence of personalized feeds. Undisclosed variation turns the form into a more concentrated version of the thing it could have exposed. Disclosure keeps it on the right side of that line. With it, a selective screening is a public event in which everyone knows that everyone is seeing the story differently. Without it, it is a private message delivered in a public room.
